% SecureGate -- Notebook 02 Scientific Report
% Designed for direct use in Overleaf.
\documentclass[10pt,a4paper]{article}

\usepackage[margin=1.8cm]{geometry}
\usepackage{amsmath,amssymb,booktabs}
\usepackage{microtype}
\usepackage{hyperref}
\usepackage{enumitem}

\setlength{\parindent}{0pt}
\setlength{\parskip}{4pt}
\setlist{nosep,leftmargin=1.4em}

\title{\textbf{SecureGate: Gate Interface and Initial Gate-Fingerprint Analysis}}
\author{Marya Alturki}
\date{\today}

\begin{document}
\maketitle

\textbf{Purpose.}
This report documents the work implemented in \texttt{02\_gate\_interface.ipynb}.
The notebook establishes the logical-gate interface used by SecureGate, audits the
available surface-code backend, implements encoded logical \(X\) and \(Z\) operations,
samples their detector records, and performs an initial gate-inference test. Its
purpose is to provide a compact scientific roadmap so that the notebook can be
followed in execution order without reconstructing the motivation of each cell.

\textbf{1. Logical-gate interface.}

A common \texttt{GateOperation} representation was introduced for logical
operations. The interface stores a gate name and the logical-qubit indices on which
the gate acts. The current SecureGate gate set is
\[
\{I,X,Y,Z,H,S,CX,T\}.
\]
The implementation normalizes gate names to uppercase, checks that the gate belongs
to the supported set, enforces the expected arity (one logical qubit for
single-qubit gates and two for \(CX\)), and requires non-negative integer qubit
indices. Basic tests verified valid \(X\) and \(CX\) objects and confirmed that an
invalid one-qubit \(CX\) operation is rejected. These tests establish a consistent
front-end representation before fault-tolerant implementations are attached.

\textbf{2. Surface-code backend audit.}

The installed \texttt{surface-sim} interface was inspected rather than assuming
that every required logical gate is natively available. For the rotated CSS
surface-code mapping used in the experiments, the directly exposed logical
operations include \(X\), \(Z\), \(S\), and \(CX\). The inspected unrotated mapping
also exposes \(H\). \(Y\) and \(T\) are not directly exposed by either mapping in
the installed interface.

Accordingly, \(Y\) and \(T\) are treated as implementation tasks requiring an
explicit construction (for example, a decomposition or state-injection strategy)
rather than being assumed to be native operations. This capability audit is
important for keeping later claims tied to the actual simulator API.

\textbf{3. Distance-5 logical \(X\) implementation.}

The gate-fingerprint experiments use a distance-\(5\) rotated surface code. The
constructed layout contains 49 physical qubits, comprising 25 data qubits and
24 ancilla qubits, and represents one logical qubit.

For the logical \(X\) operator, \texttt{surface-sim} identifies the physical
support as
\[
\{D1,D6,D11,D16,D21\},
\]
which correspond to physical indices
\[
\{0,5,10,15,20\}.
\]
The built-in \(X\)-gate configuration was inspected and applied to the layout.
The generated logical-operation block explicitly contains
\[
X\;0\;5\;10\;15\;20,
\]
with the same physical error probability applied through the noise model. The
implementation check returned \texttt{True}, confirming that the generated
physical \(X\) layer matches the logical-\(X\) support specified by the layout.

The high-level logical circuit is initialized, advanced by a \texttt{TICK}, applies
\(X\), advances by a second \texttt{TICK}, and finally measures the logical qubit.
After encoding through \texttt{surface-sim}, the complete fault-tolerant circuit
contains 49 physical qubits, 73 measurements, 60 detectors, one logical observable,
and 27 temporal ticks. The additional gates visible in the encoded circuit are
part of the syndrome-extraction procedure rather than additional logical gates.

\textbf{4. Detector sampling and \(X_L\) fingerprint.}

For each experiment, the detector sampler returns a binary record for every shot.
A detector value \(D_j=1\) denotes a detector event at detector \(j\), while
\(D_j=0\) denotes no event. The experiment was sampled for
\(10{,}000\) shots, producing a
\[
10{,}000\times 60
\]
binary detector matrix.

For detector \(j\), the empirical event probability was estimated as
\[
\widehat{P}(D_j=1\mid X_L)
=
\frac{1}{N}\sum_{n=1}^{N}D_j^{(n)},
\qquad N=10{,}000.
\]
The detector coordinates provide a spacetime location \((x,y,t)\) for each
detector, allowing the empirical fingerprint to be examined both detector-by-
detector and by temporal layer.

The measured \(X_L\) detector activity is concentrated around the logical-operation
layer. The mean detector-event probability was \(0.010087\) at \(t=0\),
\(0.035963\) at \(t=1\), and \(0.020167\) at \(t=1.5\). The largest observed
detector-event probabilities were located at \(t=1\). This establishes a
time-localized detector signature associated with the logical-operation region,
but it does not by itself establish unique gate identification.

\textbf{5. Logical \(Z\) experiment and detector correspondence.}

A separate distance-\(5\) logical \(Z\) experiment was constructed using the same
layout, detector frame, and physical error probability. The logical \(Z\) support
is the corresponding string of five data qubits \(D1\)--\(D5\), giving physical
indices \(0,1,2,3,4\).

The \(X_L\) and \(Z_L\) experiments were verified to contain the same detector
count and the same ordered spacetime detector coordinates. Therefore, detector
\(j\) in one experiment corresponds to detector \(j\) in the other experiment,
which permits direct comparison of
\[
\widehat{P}(D_j=1\mid X_L)
\quad\text{and}\quad
\widehat{P}(D_j=1\mid Z_L).
\]

The mean absolute detector-level difference was
\[
\frac{1}{60}\sum_{j=1}^{60}
\left|
\widehat{P}(D_j=1\mid X_L)
-
\widehat{P}(D_j=1\mid Z_L)
\right|
=0.001635.
\]
The maximum absolute difference was \(0.00510\), occurring at detector 26,
with spacetime coordinate
\[
(x,y,t)=(2,4,1).
\]
At this detector,
\[
\widehat{P}(D_{26}=1\mid X_L)=0.04050,
\qquad
\widehat{P}(D_{26}=1\mid Z_L)=0.03540.
\]

The temporal comparison gives
\[
\begin{array}{c|ccc}
t & X_L\ \mathrm{mean} & Z_L\ \mathrm{mean} & |\Delta|\\
\hline
0.0 & 0.010087 & 0.010271 & 0.000183\\
1.0 & 0.035963 & 0.035046 & 0.000917\\
1.5 & 0.020167 & 0.020658 & 0.000492
\end{array}
\]
so the largest aggregate difference is also concentrated at the central
logical-operation layer.

\textbf{6. Initial gate-inference experiment.}

The complete 60-detector record from each shot was used as the feature vector.
The \(10{,}000\) \(X_L\) and \(10{,}000\) \(Z_L\) shots were combined into a
balanced binary dataset. An 80/20 stratified train/test split produced
\(8{,}000\) training shots and \(2{,}000\) test shots per gate.

A Bernoulli Naive Bayes classifier was trained because detector outcomes are
binary. For an unseen detector record, the classifier estimates which gate class
is more compatible with the observed pattern. The resulting test set contained
\(4{,}000\) shots in total and gave
\[
\mathrm{Accuracy}=0.5052.
\]
The confusion matrix was
\[
\begin{array}{c|cc}
 & \widehat{X}_L & \widehat{Z}_L\\
\hline
X_L & 537 & 1463\\
Z_L & 516 & 1484
\end{array}
\]
showing a substantial bias toward predicting \(Z_L\). The approximate 95\%
confidence interval for the measured accuracy was
\[
[0.4898,\;0.5207],
\]
which contains the 0.50 chance level for two equally represented classes.

These results indicate that, under the present simulation conditions and for this
particular Bernoulli Naive Bayes baseline, the detector patterns did not provide
reliable single-shot \(X_L\)-versus-\(Z_L\) identification. This should not be
interpreted as a proof that the two gate distributions are exactly identical, or
that no inference method can ever distinguish them. Rather, small empirical
differences in detector firing probabilities did not translate into classification
performance clearly above chance.

\textbf{7. Relation to the reference work and next stage.}

The observed baseline is consistent with the reference work's treatment of Pauli
gates under an unbiased Pauli noise model: the paper reports substantial
confusion among \(X\), \(Y\), and \(Z\) gate classes and describes them as sharing
a non-distinctive fingerprint in that setting. The present experiment is narrower
and is not a direct reproduction of the paper's full setup; it tests only
\(X_L\) versus \(Z_L\) using the installed \texttt{surface-sim} model and a
Bernoulli Naive Bayes classifier.

The notebook therefore establishes a validated baseline for the next stage of
SecureGate: implementing logical gates with more distinctive fault-tolerant
signatures and evaluating whether detector spacetime patterns can support
multi-class gate inference. The distinction between native backend operations and
gates requiring explicit constructions will be preserved throughout that stage.

\textbf{Reference.}
Shukla et al., ``Anticipating Decoder Side-channel Attacks in Fault-tolerant
Quantum Computers,'' arXiv:2607.12174 (2026).

\end{document}
